\paragraph{\trAutomaticRegularity。}\label{sub:automatic-regularity}
设 $S$ 是一个带\trStriplikeEnd{}的\trConnected\trRiemann\trSurface。考虑空间%
$\widetilde{\scr{R}}$，即在某个\trCompact{}子集之外都与给定的 $j_0=I_S$ 一致的、%
$S$ 上全部\trComplexStructure{} $j$ 所构成的空间。乘积 $\widetilde{\scr{S}}=%
\widetilde{\scr{R}}\times S\to\widetilde{\scr{R}}$ 在诸\trFiber{}上带有\trTautological{}的%
\trComplexStructure，这使它成为带\trStriplikeEnd{}的\trRiemann\trSurface\trFamily。%
此外，这个\trFamily{}还带有群 $\widetilde{\scr{G}}$ 的一个自然作用，这个群由 $S$ 的%
那些\trDiffeomorphism{}组成：它们在每条\trStriplikeEnd{}的无穷附近都限制为 $s$ 方向的%
\trTranslation。
给 $S$（从而也给 $\widetilde{\scr{S}}$）配备\trLagrange\trLabel，使得对无穷远处的%
每个点 $\zeta$ 都有 $L_{\zeta,0}\pitchfork L_{\zeta,1}$。取平凡的\trPerturbationDatum{}
$K=0$、$J=I_M$（\trAlmostComplex{}结构的常值\trFamily）。取某个\trPseudoholomorphicMap{}
$u\in\scr{M}_S(\{\gamma_\zeta\})$，并把 $(j_0,u)$ 视为相应的\trParametrization{}的%
\trModuliSpace{} $\scr{M}_{\widetilde{\scr{S}}}(\{\gamma_\zeta\})$ 中的一个元素。%
$\widetilde{\scr{G}}$ 的无穷小作用定义了一个\trLinear{}映射 $d^0$，从李代数%
$L\widetilde{\scr{G}}$ 映到 \trZariski\trTangentSpace{}
$T^{\mathrm{Zar}}\scr{M}_{\widetilde{\scr{S}}}(\{\gamma_\zeta\})_{(j_0,u)}$
（我们说“\trZariski\trTangentSpace”，是因为我们没有作任何\trRegularity{}假设）。%
另一方面，这个 \trZariski\trTangentSpace{}按定义正是扩张\trLinearization{}算子%
$d^1=D_{\widetilde{\scr{S}},j_0,u}$ 的\trKernel。因此可以把这两个算子想成一个%
\trComplex{}中的\trBoundaryOperator
\begin{equation}\label{eq:automatic-complex}
0\longrightarrow L\widetilde{\scr{G}}\xrightarrow{\ d^0\ }
T\widetilde{\scr{R}}_{j_0}\times(T\scr{B}_S)_u\xrightarrow{\ d^1\ }
(\scr{E}_S)_u\longrightarrow 0.
\end{equation}
我们对 $\widetilde{\scr{R}}$ 与 $\widetilde{\scr{G}}$ 作为无穷维空间的处理一直相当不正式，%
但此刻我们想至少在\trTangentSpace{}的层面上把事情说得更精确些，办法是选取合适的%
\trHilbertSpace{}\trCompletion。如前所定义，$\widetilde{\scr{G}}$ 的李代数本应由 $S$ 上%
与边界相切的\trSmooth{}\trVectorField{} $Z$ 组成，它们在每条\trStriplikeEnd{}的无穷附近%
满足 $\epsilon_\zeta^*Z=c_\zeta\partial_s$，其中 $c_\zeta\in\RR$ 为某常数。我们取而代之的%
是它的\trWeighted\trSobolevSpace{}\trCompletion，其中 $(k,p)=(2,2)$、$\epsilon>0$ 充分小，%
即按 \eqref{eq:enlarged-domain} 的记号是 $W^{2,2;\epsilon}(S,TS,T(\partial S))^{+}$。%
类似地，$\widetilde{\scr{R}}$ 的\trTangentSpace{}天真地说应是 $TS$ 的与%
\trComplexStructure{}反交换的\trCompactlySupported{}\trEndomorphism{}所成的空间。%
我们把这些解释为（复的）\trTensorProduct\trLineBundle{} $\Omega^{0,1}_S\otimes TS$ 的%
\trSmooth{}\trSection，并\trCompletion{}到 $W^{1,2;\epsilon}(\Omega^{0,1}_S\otimes TS)$。%
最后，对\trLinearization{}算子 $D_{S,u}$——即 $d^1$ 的第二个分量——我们也使用同一套%
Sobolev \trTopology。\eqref{eq:automatic-complex} 的\trCompletion{}便形如
\begin{equation}\label{eq:fredholm-complex-completion}
\begin{gathered}
W^{2,2;\epsilon}(S,TS,T(\partial S))^{+}\\
\big\downarrow\;\, d^{0}\\
W^{1,2;\epsilon}(S,\Omega^{0,1}_S\otimes TS)\oplus W^{1,2;\epsilon}(S,E,F)\\
\big\downarrow\;\, d^{1}\\
L^{2;\epsilon}(S,\Omega^{0,1}_S\otimes E)
\end{gathered}
\end{equation}
其中 $E=u^*TM$，$F\subset E|\partial S$ 由相关的\trLagrange\trSubmanifold{}的%
\trTangentSpace{}给出，而
\begin{equation*}
\begin{aligned}
d^0(Z)&=(2I_S\circ\bar{\partial}Z,\ Du\circ Z),\\
d^1(Y,X)&=\tfrac{1}{2}I_M\circ Du\circ Y+D_{S,u}(X).
\end{aligned}
\end{equation*}
由于我们是特设地选取 Sobolev 空间的，首先应当验证这些确实是\trBounded{}算子，%
不过这很直接。实际上，还有更强的事实：$TS$ 上带\trBoundaryCondition{} $T(\partial S)$%
的标准\trCauchyRiemann{}算子 $\bar{\partial}$ 是\trAdmissible{}的，从而在\trWeighted%
\trSobolevSpace{}上成为\trFredholm{}的，讨论见 \cref{sub:kodaira-spencer}。$D_{S,u}$%
本身是一个\trNondegenerate{}的 $\bar{\partial}_\nabla$-算子，因而也是\trFredholm{}的。%
其余分量的阶为零，故是\trCompact{}的。结论是：\eqref{eq:fredholm-complex-completion}%
是一个\trFredholm\trComplex。
任取一个\trFamily{} $\scr{S}\to\scr{R}$（这一次，底空间是有限维的），使它是 $S=\scr{S}_{r_0}$%
在 \cref{sub:kodaira-spencer} 意义下的无穷小\trMiniversal{}\trDeformation。令 $\scr{G}$%
为 $S$ 的\trHolomorphic\trAutomorphism{}群。
\begin{lemma}\label{thm:fredholm-complex-cohomology}
\eqref{eq:fredholm-complex-completion} 的各\trCohomology{}群有如下解释：$H^0$ 与 $H^1$%
分别是映射 $L\scr{G}\to\ker(D_{\scr{S},r_0,u})$ 的\trKernel{}与\trCokernel，而%
$H^2=\mathrm{coker}(D_{\scr{S},r_0,u})$。
\end{lemma}
\begin{proof}
非形式地说，想法是：$\scr{R}$ 是 $\widetilde{\scr{G}}$ 在 $\widetilde{\scr{R}}$ 上的作用%
的一个\textbf{局部切片}（local slice）。为把这一想法变成严格的论证，取 $\scr{S}$ 的一个在诸\trStriplikeEnd{}上%
为常值的局部\trTrivialization。这给出 $\scr{R}$ 到 $\widetilde{\scr{R}}$ 中的一个局部%
\trEmbedding，从而给出二者\trTangentSpace{}之间的一个映射。于是我们得到一个诸列%
\trExact{}的交换图
\begin{equation}\label{eq:slice-commutative-diagram}
\begin{tikzcd}[column sep=1.05cm,row sep=1.05cm]
0\arrow[r]\arrow[d] &
T\scr{R}_{r_0}\times(T\scr{B}_S)_u\arrow[d]\arrow[r,"D_{\scr{S},r_0,u}"] &
(\scr{E}_S)_u\arrow[d,"="] \\
L\widetilde{\scr{G}}\arrow[r,"d^0"]\arrow[d,"="] &
T\widetilde{\scr{R}}_{j_0}\times(T\scr{B}_S)_u\arrow[d]\arrow[r,"d^1"] &
(\scr{E}_S)_u\arrow[d] \\
L\widetilde{\scr{G}}\arrow[r] &
T\widetilde{\scr{R}}_{j_0}/T\scr{R}_{r_0}\arrow[r] &
0.
\end{tikzcd}
\end{equation}
按定义，$T\scr{R}_{r_0}$ \trIsomorphic{}地\trProjection{}到%
\eqref{eq:dbar-plus-fredholm} 的\trCokernel{}上，这意味着%
\eqref{eq:slice-commutative-diagram} 的底行是一个 $H^0=L\widetilde{\scr{G}}$、$H^1=0$%
的\trComplex。\trCohomology{}中的\trLongExactSequence{}便给出所要的陈述。
\end{proof}
\begin{lemma}\label{thm:automatic-regularity-surface}
若 $M$ 是二维的且 $u$ 不是常值映射，则对任何无穷小\trMiniversal{}\trFamily{}
$\scr{S}\to\scr{R}$，$D_{\scr{S},r_0,u}$ 都是满的。
\end{lemma}
\begin{proof}
$u$ 是\trRiemann\trSurface{}之间的\trHolomorphic{}映射且非常值，故%
$Du\colon TS_z\to TM_{u(z)}$ 除在有限多个点之外都是\trIsomorphism。由此可见，$d^1$ 的%
第一个分量有\trDense{}的像。由第二个分量的 \trFredholm{}性质，这意味着 $d^1$ 是满的。%
然后诉诸\cref{thm:fredholm-complex-cohomology} 的最后一部分即可。
\end{proof}
我们把\cref{thm:automatic-regularity-surface}称作一条“\trAutomaticRegularity”陈述，%
因为它对任何 $I_M$ 都成立，并且不需要对\trPseudoholomorphicMap{}方程作任何%
\trPerturbation。
\begin{remark}\label{thm:sheaf-theoretic-interpretation}
惯以代数几何眼光看问题的读者，或许会觉得下面这种层论式的解读不无帮助。为简单起见，%
我们转到类似的情形：\trHolomorphic{}映射 $u\colon S\to M$，其中 $S$ 是闭%
\trRiemann\trSurface，$M$ 是\trComplex\trManifold。此时%
\eqref{eq:fredholm-complex-completion} 变为
\begin{equation*}
W^{2,2}(S,TS)\longrightarrow
W^{1,2}(S,\Omega^{0,1}_S\otimes TS)\oplus W^{1,2}(S,u^*TM)\longrightarrow
L^2(S,\Omega^{0,1}_S\otimes u^*TM),
\end{equation*}
而它事实上计算的是凝聚层的两步\trComplex{} $\{Du\colon TS\to u^*TM\}$ 的超%
\trCohomology。若 $u$ 不是常值映射，则 $Du$ 是单的（作为层映射），因此我们实际计算的%
是商 $u^*TM/TS$ 的\trCohomology（\trSuspension{}一次）。最后，若 $M$ 是一个\trSurface，%
这个商便是挠层，从而它的一次\trCohomology{}为零。
\end{remark}
